\clearpage\fancyhead[L]{\small GAUSS AN BODE.}\hypertarget{native-392-lower}{}\label{native-392-lower}
\begin{center}\textbf{GAUSS AN BODE.}\par\vspace{4pt}\hrule\vspace{4pt}\fontsize{10}{12}\selectfont Astronomisches Jahrbuch für 1818. Seite 229..232. Berlin 1815.\par\vspace{4pt}\hrule\vspace{4pt}\end{center}

\begin{flushright}
Göttingen den 9. Aug. 1815.
\end{flushright}

\par\indent Mit Vergnügen theile ich Ihnen hier meine sämmtlichen Cometenbeobachtungen mit; das so höchst
ungünstige Wetter und der Mondschein werden, fürchte ich, es wohl hindern, ihnen noch neue
beizufügen, obgleich der Comet sonst noch einige Wochen ganz gut würde beobachtet werden können.

\SourceLeaf{393}\begin{center}\fontsize{10}{12}\selectfont\setlength{\tabcolsep}{5pt}\renewcommand{\arraystretch}{1.15}\begin{tabular}{l|c|c|c}
1815 & \shortstack{Mittlere Zeit\\in Göttingen} & \shortstack{Scheinb. gerade\\Aufsteigung} & \shortstack{Scheinb. Nordl.\\Abweichung}\\\hline
März 20 & $10^{\mathrm u}\;33^{\mathrm m}\;6^{\mathrm s}$ & $54^\circ\;7'\;1''$ & $39^\circ\;7'\;47''$\\
21 & $10\;11\;37$ & $54\;33\;21$ & $39\;36\;57$\\
25 & $9\;46\;38$ & $56\;28\;50$ & $41\;38\;5$\\
30 & $9\;50\;57$ & $59\;13\;2$ & $44\;10\;27$\\
April 2 & $9\;11\;55$ & $61\;2\;27$ & $45\;39\;57$\\
Juni 12 & $10\;53\;31$ & $165\;3\;3$ & $52\;8\;25$\\
30 & $10\;49\;12$ & $184\;39\;24$ & $39\;22\;5$\\
Juli 13 & $10\;57\;20$ & $194\;32\;52$ & $29\;53\;54$\\
27 & $10\;4\;59$ & $203\;0\;56$ & $20\;37\quad:$\\
29 & $10\;13\;58$ & $204\;6\;3$ & $19\;24\;13$\\
Aug. 4 & $10\;14\;7$ & $207\;14\;8$ & $15\;53\;49$\\
\end{tabular}\end{center}

\par\indent Am 13. und 27. Juli so wie am 4. August wurden Sterne aus \textsc{Piazzis} Catalog (XII, 268; 1 Boot.;
XIII, 264. 265. 284) verglichen; an allen übrigen Tagen Sterne aus der Histoire Céleste und den
Mem. 1790. Die Beobachtungen sind alle am Kreismikrometer des 10füssigen \textsc{Herschel}schen Teleskops
angestellt, ausgenommen die vom 12. Juni, wo der Comet am \textsc{Frauenhofer}schen auf einer
vortrefflichen parallatischen Maschine aufgestellten Heliometer mit 219 Ursae majoris \textsc{Bode} (aus
den Mem. 1790 berechnet) verglichen wurde vermittelst der Distanz und des unmittelbar vom
Instrument angegebenen Positionswinkels. So vortreffliche Resultate indessen auch diese Methode
bei hellern Cometen geben mag, so ziehe ich doch in gegenwärtigem Falle (wo der lichtschwache
Comet vor dem mehr als 5mal hellern Sterne ganz verlosch, und nur erkannt werden konnte, wenn die
eine Objectivhälfte grösstentheils bedeckt war) gute Kreismikrometerbeobachtungen vor.

\Needspace{23\baselineskip}\par\indent Meine schon zu Ende März berechneten parabolischen Elemente glaube ich Ihnen schon früher
mitgetheilt zu haben. Die Abweichung derselben von der Bewegung des Cometen in den folgenden
Monaten zeigt sich beträchtlicher, als man es sonst in dergleichen Fällen gewohnt ist. Gleich
nachdem ich die Beobachtung vom 12. Juni erhalten hatte, berechnete ich aus dieser und zwei
Beobachtungen von \textsc{Olbers} vom 6. März und 25. April neue parabolische Elemente, die ich hieher
setze:

\Needspace{12\baselineskip}\begin{center}\fontsize{10}{12}\selectfont Durchgang durch die Sonnennähe 1815 April 25.\par $11^{\mathrm u}\;41^{\mathrm m}\;19^{\mathrm s}$ M. Z. in Göttingen\end{center}
\begin{center}\fontsize{9}{11}\selectfont\setlength{\tabcolsep}{2pt}\renewcommand{\arraystretch}{1.15}\begin{tabular}{lcr}
Länge der Sonnennähe & $.\;.\;.\;.\;.\;.\;.\;.\;.\;.\;.\;.$ & $147^\circ\;35'\;55''$\\
Länge des aufsteigenden Knoten & $.\;.\;.\;.\;.\;.\;.\;.$ & $82\;43\;6$\\
Neigung der Bahn & $.\;.\;.\;.\;.\;.\;.\;.\;.\;.\;.\;.\;.$ & $44\;43\;13$\\
Kleinster Abstand von der Sonne & $.\;.\;.\;.\;.\;.\;.\;.$ & $1.23024$\\
\end{tabular}\end{center}\begin{center}Bewegung rechtläufig.\end{center}

\par\indent Diese Elemnnte 
stellen zwar die drei zum Grunde gelegten Beobachtungen vollständig und fast genau
dar, allein demungeachtet weichen sie von den zwischenliegenden Beobachtungen beträchtlich ab (in
der geraden Aufsteigung bis auf 11 Minuten). Eine solche Erscheinung kann oft vorkommen, dass
drei vollständige Beobachtungen zur Bestimmung des Kegelschnitts nicht ausreichen, worüber ich in
den letzten Artikeln des ersten Abschnitts des zweiten Buches der Theoria Motus C. C.
umständlicher gesprochen habe. Da indessen durch die Abweichung der parabolischen Elemente von den
andern Beobachtungen entschieden war, dass die Bahn stark von der Parabel verschieden sein müsse,
so schien es mir interessant, eine vorläufige hypothesenfreie Berechnung der Bahn zu machen. Ich
legte dabei 5 Beobachtungen zu Grunde, die durch die herausgebrachte Ellipse bis auf Kleinigkeiten
dargestellt wurden. Auch haben diese (um die Mitte Juni berechneten und im 105. Stück der
hiesigen Gel. Anz. [S. 385 d. B.] abgedruckten) Elemente bis auf meine letzten Beobachtungen
fortwährend gute Uebereinstimmung
\SourceLeaf{394}behalten. Hr. \textsc{Nicolai} hat seitdem dieselben bereits einmal durch schärfere Rechnung und Benutzung
mehrerer Beobachtungen verbessert, und nur kleine Abänderungen erhalten. Mit Vergnügen sehe ich
auch, dass unser vortrefflicher \textsc{Bessel} sehr nahe übereinstimmende Resultate erhalten hat. Die weitere
Ausfeilung werde ich diesen geschickten Händen überlassen. Die Resultate werden ein kostbares Vermächtniss für unsere Enkel sein, die die Wiederkunft dieses höchst merkwürdigen Cometen feiern werden.

\Needspace{14\baselineskip}\begin{center}
\emph{Elliptische Elemente des Cometen von} 1815.
\end{center}

\begin{center}\fontsize{10}{12}\selectfont Durchgang durch die Sonnennähe April 26.\par $0^{\mathrm u}\;21^{\mathrm m}\;25^{\mathrm s}$ in Göttingen.\end{center}
\begin{center}\fontsize{9}{11}\selectfont\setlength{\tabcolsep}{2pt}\renewcommand{\arraystretch}{1.15}\begin{tabular}{lcr}
Länge der Sonnennähe & $.\;.\;.\;.\;.\;.\;.\;.\;.\;.\;.\;.$ & $148^\circ\;58'\;48''$\\
Abstand von der Sonnennähe & $.\;.\;.\;.\;.\;.\;.\;.\;.$ & $1.21349$\\
Aufsteigender Knoten & $.\;.\;.\;.\;.\;.\;.\;.\;.\;.\;.\;.$ & $83^\circ\;26'\;21''$\\
Neigung der Bahn & $.\;.\;.\;.\;.\;.\;.\;.\;.\;.\;.\;.\;.$ & $44\;30\;43$\\
Excentricität & $.\;.\;.\;.\;.\;.\;.\;.\;.\;.\;.\;.\;.\;.\;.$ & $0.933149$\\
\end{tabular}\end{center}

\par\indent Hieraus folgt noch Abstand in der Sonnenferne $35.0911$ Umlaufszeit $77\tfrac12$ Jahre. Diese Umlaufszeit
kann nur um wenige Jahre ungewiss sein, und wahrscheinlich ist sie noch etwas zu gross.

\Needspace{9\baselineskip}\par\indent Von der Juno, welche in der diesjährigen Opposition nur die 10. Grösse hatte, habe ich folgende
drei Beobachtungen gemacht, vielleicht die einzigen, die überhaupt diesmal angestellt sind.

\begin{center}\fontsize{9}{11}\selectfont\setlength{\tabcolsep}{2pt}\renewcommand{\arraystretch}{1.15}\begin{tabular}{llrccc}
1815 & März & 1. & $10^{\mathrm u}\;22^{\mathrm m}\;0^{\mathrm s}$ & $\mathrm{AR}=197^\circ\;7'\;56''8$ & $\mathrm{Decl.}=2^\circ\;19'\;3''0\;\mathrm{S.}$\\
 & & 29. & $10\;8\;57$ & $192\;31\;51.8$ & $1\;32\;38.1\;\mathrm{N.}$\\
 & April & 8. & $9\;27\;43$ & $190\;34\;51.9$ & $2\;53\;24.2$ ---\\
\end{tabular}\end{center}

\par\indent In Rücksicht des \textsc{Frauenhofer}schen Heliometers, wozu ich in dem letzten Frühjahr das überaus
schöne parallatische Stativ erhalten habe, bemerke ich, dass eine etwas ausführlichere Beschreibung
ohne eine genaue Zeichnung nicht wohl verständlich sein würde. Ich bin überzeugt, dass, wenn dem
Instrument eine sehr solide Aufstellung gegeben werden könnte, man bei den Bewegungen der parallatischen Maschine immer auf $\tfrac12$ Bogenminute sicher sein würde (die Verniers geben sogar $10''$ an); allein
dies erlaubt mein gegenwärtiges Local nicht. Ich muss das Instrument bald in diese bald in jene Thür
stellen; inzwischen habe ich Mittel gefunden, überall dem Instrument in sehr kurzer Zeit sogleich bis
auf 1 oder höchstens $2'$ seine richtige Stellung zu geben, und finde dann z.\,B. bei günstiger Luft den
Polarstern, den Merkur etc. bei Tage mit leichter Mühe.
